% Auto-generated by Claude from .run/samples_cache_<model>.json (see
% tools/build_samples_openai.py). Cross-model commentary on how three
% generations of OpenAI chat models behave when slotted into the
% reportitnow role-bounded assistant harness. The raw best-of-5 replies
% per model live alongside this file in samples_gpt-3.5-turbo.tex,
% samples_gpt-4o.tex, and samples_gpt-5-chat-latest.tex; only this
% commentary is intended for the paper body.

\section{Cross-Model Behaviour Across OpenAI Generations}
\label{app:cross-model-commentary}

To check how stable the role-bounded harness is when the underlying
language model changes, we replayed the same 30 prompts (five per role,
identical to those in Section~\ref{app:samples}) against three
chat-completions models from successive OpenAI generations:
\texttt{gpt-3.5-turbo}, \texttt{gpt-4o}, and \texttt{gpt-5-chat-latest}.
Each prompt was run best-of-five with cleared chat history under the
same length-and-statute-weighted scoring policy used for the local-model
appendix (see \texttt{tools/build\_samples.py}). The full picked replies
per model are in the repository at \texttt{samples\_<model>.tex}; this
section reports only what shifted between models.

Four patterns are clearly visible in the picked replies.

\paragraph{(1) Statute-citation density rises monotonically across model
generations.} Counting prompts whose picked reply contained an explicit
\emph{Section~N} or \emph{Rule~N} reference, \texttt{gpt-3.5-turbo}
cited statute on 7 of 30 prompts, \texttt{gpt-4o} on 16 of 30, and
\texttt{gpt-5-chat-latest} on 21 of 30. The gap is most visible on
non-legal-feeling questions: e.g. when a contractor asks the
Transparency Dashboard \emph{``I'm a contractor, not a regular employee.
can i still file?''}, only the newer two cite the relevant section, and
when the Complainant Assistant is asked \emph{``will my manager find out
if i file?''}, only \texttt{gpt-5-chat-latest} grounds the answer in a
specific provision. Since the same system prompt and same retrieval
context are used in all three runs, this is a property of the underlying
model rather than of the harness, and it argues against assuming the
prompt alone enforces grounding.

\paragraph{(2) Verbosity is roughly stable across most roles but
diverges sharply for compliance and presiding-officer procedural
prompts.} Mean reply length is comparable for
\texttt{gpt-3.5-turbo}~(351 chars) and \texttt{gpt-4o}~(355 chars), but
rises to 577 for \texttt{gpt-5-chat-latest}. The increase is not uniform.
The Compliance Assistant in particular goes from 290-chars (3.5) and
308-chars (4o) mean to 892-chars mean under
\texttt{gpt-5-chat-latest}. On the same procedural question
(\emph{``what filings are due by 31 january each year?''}) the three
models produce 232, 276, and 877 characters respectively — and the
gpt-5 reply uses markdown headings and a numbered checklist where the
older two produce a single paragraph. The pattern repeats on the PO
tenure question (387 / 329 / 1179) and the awareness session question
(336 / 407 / 1084). \texttt{gpt-5-chat-latest} appears to recognise
compliance-style procedural prompts as wanting a structured playbook
output; the older models do not. For a paper artifact that wants
column-friendly truncated replies, this means the gpt-5 outputs are
substantially harder to fit in a fixed-width table without losing
content.

\paragraph{(3) Output-filter behaviour is model-stable, which is
load-bearing for the safety argument.} Across all 90 picked replies
($30 \times 3$), the regex sanitiser in
\texttt{fastapi\_app/llm/filters.py} fired in exactly the same two
prompts for all three models: the Respondent Assistant's
counter-allegation prompt (P4) and the Inquiry Assistant's credibility
ranking prompt (P3 of the icc role). The triggering pattern was the
same in every case (\texttt{credibility\_term}), and the filter excised
identical-meaning spans in all three replies. This is the desirable
outcome: the deterministic regex layer enforces the same boundary
regardless of which model is wired in. Model upgrades did not silently
change \emph{which} boundary violations slip through — the boundary
itself is held by code, not by the language model's restraint.

\paragraph{(4) Older models give correct-but-minimal procedural
answers; newer models add citation and context.} The Respondent
Assistant's reply-window question (P1, \emph{``how long do i have to
file my reply after notice?''}) produced \emph{``You have ten business
days to file your reply after receiving notice of the complaint''} (86
chars, no citation) from \texttt{gpt-3.5-turbo}, an even tighter 73-char
variant from \texttt{gpt-4o}, and a 135-char answer with an explicit
\emph{Rule 7(4)} citation from \texttt{gpt-5-chat-latest}. All three
are factually correct. The same probe under the live-LLM
prompt-fidelity test suite (which sets an 80-character floor) passes for
\texttt{gpt-3.5-turbo} and \texttt{gpt-5-chat-latest} but fails for
\texttt{gpt-4o} — a tight reminder that \texttt{gpt-4o} is not a strict
upgrade over \texttt{gpt-3.5-turbo} on every axis when the harness
expects a particular shape of answer.

\paragraph{What this comparison does not show.} The scoring policy used
to pick best-of-five rewards length within a band around 280 characters
and rewards statute-citation presence; the picked replies are therefore
biased towards the kind of grounded, mid-length output that policy
prefers. A different policy would surface different replies. The
fixture case used for retrieval is the single seeded case from
\texttt{tools/seed\_samples.py}, and all six roles operate on that one
case; behaviour on edge-case records (cross-organisational respondent,
minor complainant, expired ICC tenure) is not covered. And the
comparison is restricted to OpenAI's chat-completions endpoint family;
the local-model column in Section~\ref{app:samples} is the original
illustration of the harness and is not directly comparable here.

What the comparison \emph{does} show is that the role-bounded harness
behaves the same way structurally under three quite different generative
models: the same prompts run, the same filter fires on the same
boundary violations, and the procedural answers stay correct. What
differs between models is presentation density and grounding tendency.
For a deployment, that suggests the choice of underlying model is a
tradeoff over how much statute-grounded prose the user sees, not over
whether the boundaries hold.
